\section{Multi-Step Proximal Policy Improvement}
\label{sec:method}

At each delayed actor update, MPI maintains $K$ independently initialized,
persistent deterministic actors with separate optimizer states. The first
actor takes one TD3+BC update toward the sampled dataset actions; each later
actor takes one update toward the stop-gradient output of its predecessor on
the same minibatch. The online critic is fixed within this chain, and its first
member is $\widehat Q$. A Polyak copy of the first actor supplies TD-target
actions, and the final actor is evaluated. We first derive the proximal
operator underlying these updates and then formalize their composition.
\subsection{Offline actor updates as proximal improvement}

Let $\rho_{\mathcal D}$ be the empirical state distribution and consider a
deterministic actor $\mu:\mathcal S\rightarrow\mathcal A\subset\mathbb R^d$.
At one training iteration, the critic parameters are held fixed while the actor
chain is updated. We write the current critic as $\widehat Q$ and normalize its
global scale by
\begin{equation}
\begin{aligned}
 \bar Q_k(s,a)&=\frac{\widehat Q(s,a)}{C_k},\\
 C_k&=\operatorname{sg}\!\left(
 \mathbb E_s|\widehat Q(s,\widetilde\mu_k(s))|+\epsilon\right).
\end{aligned}
\label{eq:qnorm}
\end{equation}
where $\widetilde\mu_k$ is the actor at the beginning of the $k$th update and
$\operatorname{sg}$ denotes stop-gradient. For fixed positive $C$, normalized
and raw-critic ascent have the same integral curves after a constant change of
time. Recomputing a positive global scalar $C_k$ preserves the instantaneous
raw-critic ascent direction and gives a trajectory-dependent time
reparameterization in the infinitesimal limit; at finite steps it changes the
effective step schedule. All finite-step statements condition on realized
$C_k$.

For deterministic policies and a state distribution $\rho$, define the
action-averaged Wasserstein metric
\begin{equation}
 \mathsf d_{d,\rho}^2(\mu,\nu)
 :=\mathbb E_{s\sim\rho}
 \left[\frac{1}{d}\|\mu(s)-\nu(s)\|_2^2\right].
 \label{eq:mean_metric}
\end{equation}
It is $d^{-1}$ times the state-averaged squared $W_2$ distance between
state-conditioned Dirac policies. We abbreviate
$\mathsf d_d=\mathsf d_{d,\rho_{\mathcal D}}$ on the actor-training state
distribution.
For diagonal-Gaussian policies, the statewise counterpart is
$W_2^2=\|m-m^{\rm ref}\|_2^2+\|\sigma-\sigma^{\rm ref}\|_2^2$, so the
same construction acts on mean and standard-deviation coordinates. The
deterministic case is its zero-variance limit; the formal statements and
fixed-budget mechanism study below use deterministic actors.

The population counterpart of squared behavior cloning also admits a
deterministic reference. Let $P_{\!B}$ be the data-generating behavior
distribution, let $\rho_B$ be its state marginal, and define
$b(s)=\mathbb E_{P_{\!B}}[a\mid s]$. Then
\begin{equation}
 \mathbb E_{(s,a)\sim P_{\!B}}\frac{1}{d}\|\mu(s)-a\|_2^2
 =\mathsf d_{d,\rho_B}^2(\mu,b)+\mathrm{const},
 \label{eq:bc_decomposition}
\end{equation}
where the constant is the conditional action variance and does not depend on
$\mu$. The empirical actor loss is a Monte Carlo estimate of the left-hand
side, so samplewise squared anchoring targets the conditional-mean behavior map
at the population level.

Given a reference policy $\nu$ and step $h>0$, define single-step proximal
improvement (SPI) by
\begin{equation}
\operatorname{SPI}_{h}^{\widehat Q}(\nu)
\in\argmin_{\mu\in\mathcal M}
\left\{-\mathbb E_s\bar Q_k(s,\mu(s))
+\frac{1}{2h}\mathsf d_d^2(\mu,\nu)\right\}.
\label{eq:spi}
\end{equation}
Multiplying by $2h$ gives the practical actor objective
\begin{equation}
\mathcal L_{h,k}(\mu\mid\nu)
=-2h\,\mathbb E_s\bar Q_k(s,\mu(s))
+\mathsf d_d^2(\mu,\nu).
\label{eq:actorloss}
\end{equation}

\begin{proposition}[Behavior anchoring is Wasserstein proximal improvement]
\label{prop:spi_equiv}
For deterministic policies, the normalized anchored objective
\begin{equation*}
 -\alpha\,\mathbb E_s\bar Q(s,\mu(s))
 +\mathsf d_d^2(\mu,\nu)
\end{equation*}
is Eq.~\eqref{eq:spi}, up to multiplication by a positive constant, with
$h=\alpha/2$. Equivalently, the raw-critic objective
$-\lambda\mathbb E_s\widehat Q(s,\mu(s))+\mathsf d_d^2(\mu,\nu)$ has
$h=\lambda C/2$ when $C$ is fixed. The population counterpart of squared
dataset anchoring uses $\nu=b$ by Eq.~\eqref{eq:bc_decomposition}.
\end{proposition}

The proposition identifies an update operator relative to the \emph{estimated}
critic; true-return improvement lies outside the scope of this certificate.

\subsection{Re-centered composition and actor--critic dataflow}

For population analysis, let $\mu_0:=b$ denote the conditional-mean behavior
map and let $h_{1:K}$ be positive step sizes. MPI constructs
\begin{equation}
 \mu_k\in\operatorname{SPI}_{h_k}^{\widehat Q}
 (\operatorname{sg}\mu_{k-1}),
 \qquad k=1,\ldots,K.
 \label{eq:mpi}
\end{equation}
The implemented first actor $\mu_1$ uses sampled dataset actions as its
reference. Later actors $\mu_2,\ldots,\mu_K$ use the preceding online actor
under the same current critic. Only the Polyak copy of $\mu_1$ enters the TD
target, while $\mu_K$ supplies the evaluated action.

\begin{figure}[t]
\centering
\begin{tikzpicture}[x=1cm,y=1cm,>=stealth, every node/.style={font=\small}]
  \node[draw,rounded corners,minimum width=1.05cm,minimum height=.52cm] (p0) at (0,0) {$\mu_0$};
  \node[draw,rounded corners,minimum width=1.05cm,minimum height=.52cm] (p1) at (1.65,0) {$\mu_1$};
  \node[draw,rounded corners,minimum width=1.05cm,minimum height=.52cm] (p2) at (3.30,0) {$\mu_2$};
  \node at (4.35,0) {$\cdots$};
  \node[draw,rounded corners,minimum width=1.05cm,minimum height=.52cm] (pk) at (5.45,0) {$\mu_K$};
  \draw[->] (p0)--node[above] {$h_1$}(p1);
  \draw[->] (p1)--node[above] {$h_2$}(p2);
  \draw[->] (p2)--(4.05,0);
  \draw[->] (4.65,0)--node[above] {$h_K$}(pk);
  \node[draw,rounded corners,minimum width=1.55cm,minimum height=.52cm] (critic) at (1.65,-1.08) {TD target};
  \draw[->,thick] (p1)--node[right,font=\scriptsize] {target copy}(critic);
  \node[below=0.05cm of p1,font=\scriptsize] {$D_{\rm critic}$};
  \node[below=0.05cm of pk,font=\scriptsize] {$D_{\rm final}$};
  \draw[decorate,decoration={brace,mirror,amplitude=4pt},yshift=-0.48cm]
  (p0.south west)--node[below=5pt,font=\scriptsize] {final-policy reach}(pk.south east);
\end{tikzpicture}
\caption{MPI re-centers each actor update. A Polyak copy of $\mu_1$ supplies
TD-target actions, while later actors increase or reshape final-policy reach
without entering the target policy.}
\label{fig:architecture}
\end{figure}

Our fixed-budget schedule sets $h_k=T/K$ and holds the total nominal budget
$T$ fixed. Increasing $K$ jointly changes local step size, re-centering
frequency, actor compute, and the target/deployment dataflow; the experiment
therefore estimates the effect of this implemented refinement bundle.

On the next-state distribution $\rho_{\mathcal D}'$ used by the TD target,
define the action-averaged population displacements
\begin{equation}
\begin{aligned}
D_{\rm critic}&:=\mathsf d_{d,\rho_{\mathcal D}'}^2(\mu_1,\mu_0),\\
D_{\rm final}&:=\mathsf d_{d,\rho_{\mathcal D}'}^2(\mu_K,\mu_0).
\end{aligned}
\label{eq:displacements}
\end{equation}
A one-step actor has $D_{\rm critic}=D_{\rm final}$; MPI assigns the two
population quantities to distinct policy branches. On a fixed current-state
batch, we define the sample-anchored online-actor proxies
\begin{equation}
\begin{aligned}
\widehat D_{\rm critic}
&:=\mathbb E_i\frac{1}{d}\|\mu_1(s_i)-a_i\|_2^2,\\
\widehat D_{\rm final}
&:=\mathbb E_i\frac{1}{d}\|\mu_K(s_i)-a_i\|_2^2.
\end{aligned}
\label{eq:sample_displacements}
\end{equation}
The direct audit reproduces the training transition loader and retains only
rows with a nonzero bootstrap mask; timeout boundaries and terminal source
rows are excluded. On these rows, we also measure the deterministic
target-policy displacement
\begin{equation}
 \widehat D_{\rm target}:=
 \mathbb E_i\frac{1}{d}
 \|\mu_1^-(s_i')-a_{i+1}\|_2^2,
 \label{eq:direct_target_displacement}
\end{equation}
where $\mu_1^-$ is the Polyak target actor and $a_{i+1}$ is the paired next
dataset action. A smoothed counterpart uses the clipped TD3 target action with
common noise draws. Sample anchoring adds the same conditional-action-variance
floor in population comparisons; paired differences cancel this floor. These
quantities measure policy displacement rather than value-target perturbation
or critic reliability.

\subsection{Finite-step statements}

The following statements do not require a small-step Taylor expansion.

\begin{proposition}[Inexact estimated-energy margin]
\label{prop:margin}
Suppose $\mu_{k-1}$ is feasible for the $k$th actor class and the computed
actor satisfies
\begin{equation*}
\mathcal L_{h_k,k}(\mu_k\mid\mu_{k-1})
\leq
\mathcal L_{h_k,k}(\mu_{k-1}\mid\mu_{k-1})+\varepsilon_k.
\end{equation*}
Then
\begin{equation}
\begin{split}
 m_k^{\widehat Q}:={}&
 \frac{2h_k}{C_k}\mathbb E_s
 [\widehat Q(s,\mu_k(s))-\widehat Q(s,\mu_{k-1}(s))]\\
 &-\mathsf d_d^2(\mu_k,\mu_{k-1})
 \geq-\varepsilon_k.
\end{split}
\label{eq:margin}
\end{equation}
\end{proposition}

For the first sample-anchored hop, a statewise dataset action is not generally
a feasible deterministic-policy comparator; the population reference in
Eq.~\eqref{eq:bc_decomposition} or a behavior actor is required for a literal
certificate. For re-centered hops, the preceding actor is feasible by
construction.

The margin is relative to $\widehat Q$. Let a trusted objective be
$Q^\dagger$ and write $\widehat Q=Q^\dagger+e$. Then
\begin{equation}
\begin{split}
 m_k^{Q^\dagger}=m_k^{\widehat Q}
 -\frac{2h_k}{C_k}\mathbb E_s[&e(s,\mu_k(s))\\
 &-e(s,\mu_{k-1}(s))].
\end{split}
\label{eq:critic_error_margin}
\end{equation}
Thus estimated-energy descent transfers to $Q^\dagger$ only when critic error
is controlled along the realized path. Equation~\eqref{eq:margin} certifies
the learned objective rather than monotone true return.

\begin{proposition}[First-hop TD-target perturbation bound]
\label{prop:exposure}
For clarity, consider a noiseless TD target using $\mu_1$,
$y_{\mu_1}=r+\gamma\min_jQ_j^-(s',\mu_1(s'))$, and suppose each target critic
$Q_j^-(s',\cdot)$ is $L(s')$-Lipschitz. For any reference $\mu_0$,
\begin{equation}
\begin{split}
\|y_{\mu_1}-y_{\mu_0}\|_{L_2(\mathcal D)}
\leq\gamma\Big(&\mathbb E_{s'}L(s')^2\\
&\cdot\|\mu_1(s')-\mu_0(s')\|_2^2\Big)^{1/2}.
\end{split}
\label{eq:exposure_bound}
\end{equation}
If $L(s')\leq\bar L$, the right-hand side is at most
$\gamma\bar L\sqrt{dD_{\rm critic}}$. Conditional on $\mu_1$, later actors do
not enter this one-step target perturbation.
\end{proposition}

Proposition~\ref{prop:exposure} bounds a one-step value-target difference, not
the stability of the learned critic. Its architectural content is that later
refinement can alter $D_{\rm final}$ without directly entering this bound.
With a Polyak target network, $\mu_1$ is replaced by its target copy. A
same-noise comparison under TD3 target smoothing obeys the same bound because
action clipping is nonexpansive. The current-state
$\widehat D_{\rm critic}$ omits Polyak lag and target noise, whereas
$\widehat D_{\rm target}$ measures the Polyak actor directly and has a
smoothed counterpart. Neither displacement measures the action Lipschitz
factor or the resulting change in critic value.

The separation introduces a complementary evaluation-policy mismatch: critic
targets follow the $\mu_1$ branch, whereas evaluation follows $\mu_K$.
Increasing $D_{\rm final}$ at fixed $D_{\rm critic}$ can therefore make the
critic less representative of the evaluated policy. We interpret MPI as
exposing a measurable target/deployment trade-off; its empirical return
captures the net effect of target movement, evaluation mismatch, actor compute,
and the coupled optimization trajectory.

\subsection{Local Wasserstein-flow limit}
\label{sec:local_limit}

The gradient-flow interpretation is a local organizing principle. Fix a critic
and an interior action $a_0$. Let $g=\nabla_a\widehat Q(a_0)$,
$H=\nabla_a^2\widehat Q(a_0)$, and $\beta=dh/C$. Before projection, the
scale-matched linearized and exact statewise proximal updates satisfy
\begin{equation}
\Delta a_{\rm lin}=\beta g,
\qquad
\Delta a_{\rm prox}=\beta\nabla_a\widehat Q(a_0+\Delta a_{\rm prox}).
\label{eq:local_steps}
\end{equation}
When $\beta\|H\|\ll1$, the latter yields
\begin{equation}
\Delta a_{\rm prox}=(I-\beta H)^{-1}\beta g
=\beta g+\beta^2Hg+O(\beta^3),
\label{eq:resolvent}
\end{equation}
and both updates are first-order discretizations of
$\dot a=d\nabla_a\bar Q(a)$. With fixed $C$ this is a constant
time-rescaling of raw-critic gradient ascent; with a positive policy-dependent
global $C$ it follows the same direction under an adaptive time
parameterization. This expansion additionally requires small realized
displacement, inactive projection, and a sufficiently accurate actor solve.
At large budgets, our experiments diagnose the resulting nonlocal
actor--critic regime instead of extrapolating Eq.~\eqref{eq:resolvent}. The
end-to-end return sweep, even at small nominal $T$, is not a frozen-critic
action-level consistency test of this limit.

\begin{proposition}[Local fixed-budget composition]
\label{prop:local_composition}
Let a smooth one-step update in local coordinates satisfy
$\Psi_h(x)=x+h f(x)+h^2p(x)+O(h^3)$. For fixed $K$ and $T\rightarrow0$,
\begin{equation}
\begin{aligned}
\Psi_{T/K}^{K}(x_0)
={}&x_0+Tf_0\\
&+T^2\!\left[\frac{K-1}{2K}J_0f_0+\frac{1}{K}p_0\right]
+O(T^3).
\end{aligned}
\label{eq:k_comp}
\end{equation}
where $f_0=f(x_0)$ and $J_0=Df(x_0)$. Coordinate explicit Euler has
$p_0=0$, whereas exact backward Euler has $p_0=J_0f_0$. Hence both are
first-order methods; for each fixed $K$, the $T\to0$ expansion has a leading
flow-bias coefficient proportional to $1/K$.
\end{proposition}

Proposition~\ref{prop:local_composition} explains a possible fixed-critic
small-step benefit of substepping. It neither predicts the sign of true-return
change nor applies to the projected, saturated, moving-critic branches studied
at large $T$. The expansion is not a uniform $K\to\infty$ result.
